\section{Definitions and Examples}

\begin{enumerate}
\item An \textbf{$\R$-linear equation} in the variables $x_1, \dots,
  x_n$ is an equation that can be expressed in the form
  \begin{displaymath}
    a_1 x_1 + \dots + a_n x_n = b
  \end{displaymath}
  where $a_1, \dots, a_n, b \in \R$.\footnote{In this course we will
  often drop the ``$\R$'' and simply call this a linear equation.}
  The numbers $a_1, \dots, a_n$ are called \textbf{coefficients} and
  the variables $x_1, \dots, x_n$ are also called \textbf{unknowns}.

  \textit{Examples.}
  \begin{itemize}
  \item $2x - 3y + \pi z = -44.5$ is a linear equation in the
    variables $x, y, z$.\footnote{The choice of variables names is
    unimportant, and will be clear from context.}
  \item $2(x_1 + x_2) = 4x_3 - x_4$ is a linear equation in the
    variables $x_1, x_2, x_3, x_4$.\footnote{This equation can be
    \textit{expressed} in the above form by rearranging terms.}
  \item $(2x_1 + x_2)(x_3 - x_4) = 5$ is not a linear equation because
    the left-hand-side, after rearranging terms, includes
    nonlinear terms, e.g., $x_2x_3$.
  \end{itemize}
\item A \textbf{system of linear equations} in the variables $x_1,
  \dots, x_n$ is a collection of linear equations in those variables.
  We also call this a \textbf{linear system}.  Note that every
  equation in a linear system must over the \textit{same} variables.
  The general form of a system of linear equations is written as
  follows.
  \begin{equation}
    \begin{split}
      a_{11}x_1 + \dots + a_{1n} x_n &= b_1 \\
      &\vdots \\
      a_{m1}x_1 + \dots + a_{mn} x_n &= b_m
    \end{split}
    \label{gen-form-lin-sys}
  \end{equation}

  \textit{Examples.}
  \begin{itemize}
  \item The following is a system of two linear equations over the
    variables $x$ and $y$.  This system represents a pair of lines in
    the plane that intersect at a single point.
    \begin{equation}
      \begin{split}
        2x + 3y &= 4 \\
        -5x + 3y &= 0
      \end{split}
      \label{lin-sys-ex-1}
    \end{equation}
  \item The following is a system of three linear equations over four
    variables.  Note that equations are required to have nonzero
    coefficients for every variable.
    \begin{equation}
      \begin{split}
      3x_1 - 4.4x_2 &= 1 \\
      200x_1 + 200x_2 + x_3 + 200x_4 &= 200 \\
      9x_1 - x_2 - 5x_4 &= 98
      \end{split}
      \label{lin-sys-ex-2}
    \end{equation}
  \end{itemize}
\item
  We write $\R^n$ for \textbf{$n$-dimensional Euclidean space}.  In
  this chapter we write $(a_1, \dots, a_n)$ where $a_1, \dots, a_n \in
  R$ for an element of $\R^n$.  We call this a \textbf{point} in
  $\R^n$.\footnote{In the next chapter, we will identify points in
  $\R^n$ with \textit{vectors}.}
\item
  An \textbf{$\R$-matrix} is a rectangular grid of real numbers.  We
  usually assume that a matrix has at least one row and and one
  column.  We write $\R^{m \times n}$ for the set of all matrices with
  $m$ rows and $n$ columns.  We identify $\R^{1 \by 1}$ with $\R$
  (i.e., we don't make a distinction between real numbers and matrices
  with one entry).

  \textit{Examples.}
  \begin{itemize}
  \item The following is a matrix in $\R^{3 \by 2}$.
    \begin{displaymath}
      \begin{bmatrix}
        1 & 2 \\
        \pi & -4 \\
        -22.4 & 0
      \end{bmatrix}
    \end{displaymath}
  \item The following is a matrix in $\R^{1 \by 5}$.
    \begin{displaymath}
      \begin{bmatrix}
        1 & 2 & 3 & 4 & 5
      \end{bmatrix}
    \end{displaymath}
  \end{itemize}
\item The \textbf{augmented matrix} of a system of $m$ linear
  equations over $n$ variables is the matrix in $\R^{m \by (n + 1)}$
  whose entries include the coefficients and right-hand-sides of each
  linear equation in the system.  The augmented matrix of the a
  general linear system (as in \ref{gen-form-lin-sys}) is as follows.
  \begin{displaymath}
    \begin{bmatrix}
      a_{11} & \dots & a_{1n} & b_1 \\
      \vdots & \ddots & \vdots & \vdots \\
      a_{m1} & \dots & a_{mn} & b_m
    \end{bmatrix}
  \end{displaymath}

  \textit{Examples.}
  \begin{itemize}
  \item
    The following matrix in $\R^{2 \by 3}$ is the augmented matrix of
    the linear system in \ref{lin-sys-ex-1}.
    \begin{displaymath}
      \begin{bmatrix}
        2 & 3 & 4 \\
        -5 & 3 & 0
      \end{bmatrix}
    \end{displaymath}
  \item
    The following matrix in $\R^{3 \by 5}$ is the augmented matrix of
    the linear system in \ref{lin-sys-ex-2}.
    \begin{displaymath}
      \begin{bmatrix}
        3 & -4.4 & 0 & 0 & 1 \\
        200 & 200 & 1 & 200 & 200 \\
        9 & -1 & 0 & -5 & 98
      \end{bmatrix}
    \end{displaymath}
  \end{itemize}
\item The \textbf{coefficient matrix} of a system of $m$ linear
  equations over $n$ variables is the matrix in $\R^{m \by n}$ gotten
  by removing the rightmost column of its augmented matrix.  The
  coefficient matrix of a general linear system (as in
  \ref{gen-form-lin-sys}) is as follows.
  \begin{align*}
    \begin{bmatrix}
      a_{11} & \dots & a_{1n} \\
      \vdots & \ddots & \vdots \\
      a_{m1} & \dots & a_{mn}
    \end{bmatrix}
  \end{align*}

  \textit{Examples.}
  \begin{itemize}
  \item
    The following matrix in $\R^{2 \by 2}$ is the coefficient matrix of
    the linear system in \ref{lin-sys-ex-1}.
    \begin{displaymath}
      \begin{bmatrix}
        2 & 3 \\
        -5 & 3
      \end{bmatrix}
    \end{displaymath}
  \item
    The following matrix in $\R^{3 \by 5}$ is the coefficient matrix of
    the linear system in \ref{lin-sys-ex-2}.
    \begin{displaymath}
      \begin{bmatrix}
        3 & -4.4 & 0 & 0 \\
        200 & 200 & 1 & 200 \\
        9 & -1 & 0 & -5
      \end{bmatrix}
    \end{displaymath}
  \end{itemize}
\item A \textbf{solution} of a linear system over the variables $x_1,
  \dots, x_n$ is a point $(s_1, \dots, s_n) \in \R^n$ such that
  substituting each variable $x_i$ with the value $s_i$ simultaneously
  satisfies every equation in the system.  A solution of the linear
  general system in \ref{gen-form-lin-sys} is a point $(s_1, \dots, s_n) \in
  \R^n$ such that the following real number equalities hold.
  \begin{align*}
    a_{11}s_1 + \dots + a_{1n} s_n &= b_1 \\
    &\vdots \\
    a_{m1}s_1 + \dots + a_{mn} s_n &= b_m \\
  \end{align*}

  \textit{Examples.}
  \begin{itemize}
  \item The point $\left(\frac{12}{21}, \frac{20}{21}\right)$ is a
    solution to the linear system in \ref{lin-sys-ex-1} because
    \begin{displaymath}
      2\left(\frac{12}{21}\right) + 3\left(\frac{20}{21}\right)
      = \frac{24}{21} + \frac{60}{21}
      = \frac{84}{21}
      = 4
      \quad
      \text{and}
      \quad
      -5\left(\frac{12}{21}\right) + 3\left(\frac{20}{21}\right)
      = -\frac{60}{21} + \frac{60}{21}
      = 0
    \end{displaymath}
  \item
    \todo{Another example.}
  \end{itemize}
\item A linear system is \textbf{consistent} if it has a solution.
  Otherwise, it is \textbf{inconsistent}.

  \textit{Examples.}
  \begin{itemize}
  \item As noted above, the linear system in \ref{lin-sys-ex-1} is consistent.
  \item The following linear system is inconsistent.  It represents
    two parallel lines in $\R^2$.
    \begin{align*}
      2x - 3y &= 4 \\
      -4x + 6y &= 5
    \end{align*}
  \end{itemize}

\item There are three kinds of \textbf{elementary row operations}.
  \begin{enumerate}
  \item \textbf{Replacement} operations add a scalar multiple of one
    row to another row.  They are denoted
    \begin{displaymath}
      R_i \gets R_i + cR_j
    \end{displaymath}
    where $c \in \R$ and act on matrices as follows.
    \begin{displaymath}
      \begin{bmatrix}
        a_{11} & \dots & a_{1n} \\
        \vdots & \ddots & \vdots \\
        a_{i1} & \dots & a_{in} \\
        \vdots & \ddots & \vdots \\
        a_{j1} & \dots & a_{jn} \\
        \vdots & \ddots & \vdots \\
        a_{m1} & \dots & a_{mn}
      \end{bmatrix}
      \xrightarrow{R_i \gets R_i + cR_j}
      \begin{bmatrix}
        a_{11} & \dots & a_{1n} \\
        \vdots & \ddots & \vdots \\
        a_{i1} + ca_{j1} & \dots & a_{in} + ca_{jn} \\
        \vdots & \ddots & \vdots \\
        a_{j1} & \dots & a_{jn} \\
        \vdots & \ddots & \vdots \\
        a_{m1} & \dots & a_{mn}
      \end{bmatrix}
    \end{displaymath}
  \item \textbf{Scaling} operations multiply a row by a nonzero
    constant. They are denoted
    \begin{displaymath}
      R_i \gets cR_i
    \end{displaymath}
    where $c \in \R$ and $c \neq 0$ and act on matrices as follows.
    \begin{displaymath}
      \begin{bmatrix}
        a_{11} & \dots & a_{1n} \\
        \vdots & \ddots & \vdots \\
        a_{i1} & \dots & a_{in} \\
        \vdots & \ddots & \vdots \\
        a_{m1} & \dots & a_{mn}
      \end{bmatrix}
      \xrightarrow{R_i \gets cR_i}
      \begin{bmatrix}
        a_{11} & \dots & a_{1n} \\
        \vdots & \ddots & \vdots \\
        ca_{i1} & \dots & ca_{in} \\
        \vdots & \ddots & \vdots \\
        a_{m1} & \dots & a_{mn}
      \end{bmatrix}
    \end{displaymath}
  \item \textbf{Swap} operations interchange two rows.  They are
    denoted
    \begin{displaymath}
      R_i \leftrightarrow R_j
    \end{displaymath}
    and act on matrices as follows.
    \begin{displaymath}
      \begin{bmatrix}
        a_{11} & \dots & a_{1n} \\
        \vdots & \ddots & \vdots \\
        a_{i1} & \dots & a_{in} \\
        \vdots & \ddots & \vdots \\
        a_{j1} & \dots & a_{jn} \\
        \vdots & \ddots & \vdots \\
        a_{m1} & \dots & a_{mn}
      \end{bmatrix}
      \xrightarrow{R_i \leftrightarrow R_j}
      \begin{bmatrix}
        a_{11} & \dots & a_{1n} \\
        \vdots & \ddots & \vdots \\
        a_{j1} & \dots & a_{jn} \\
        \vdots & \ddots & \vdots \\
        a_{i1} & \dots & a_{in} \\
        \vdots & \ddots & \vdots \\
        a_{m1} & \dots & a_{mn}
      \end{bmatrix}
    \end{displaymath}
  \end{enumerate}

  \textit{Examples.}
  \begin{itemize}
  \item The following is an example of applying the replacement
    operation $R_2 \gets R_2 - 3R_1$.
    \begin{displaymath}
      \begin{bmatrix}
        1 & 2 & 3 \\
        4 & 5 & 6 \\
        7 & 8 & 9
      \end{bmatrix}
      \xrightarrow{R_2 \gets R_2 - 3R_1}
      \begin{bmatrix}
        1 & 2 & 3 \\
        1 & -1 & -3 \\
        7 & 8 & 9
      \end{bmatrix}
    \end{displaymath}
  \item The following is an example of applying the scaling operation $R_3 \gets (-4)R_3$.
    \begin{displaymath}
      \begin{bmatrix}
        1 & 2 & 3 \\
        4 & 5 & 6 \\
        7 & 8 & 9
      \end{bmatrix}
      \xrightarrow{R_3 \gets (-4)R_3}
      \begin{bmatrix}
        1 & 2 & 3 \\
        4 & 5 & 6 \\
        -28 & -32 & -36
      \end{bmatrix}
    \end{displaymath}
  \item The following is an example of applying the swap operation $R_1 \leftrightarrow R_2$.
    \begin{displaymath}
      \begin{bmatrix}
        1 & 2 & 3 \\
        4 & 5 & 6 \\
        7 & 8 & 9
      \end{bmatrix}
      \xrightarrow{R_1 \leftrightarrow R_2}
      \begin{bmatrix}
        4 & 5 & 6 \\
        1 & 2 & 3 \\
        7 & 8 & 9
      \end{bmatrix}
    \end{displaymath}
  \end{itemize}
\item Matrices $A, B \in \R^{m \times n}$ are \textbf{row equivalent},
  written $A \sim B$, if there is a sequence of elementary row
  operations that transform $A$ into $B$.

  \textit{Facts.}
  \begin{itemize}
  \item
    \textit{Row equivalence is, in fact, an equivalence.}
    \begin{itemize}
    \item (reflexivity) $A \sim A$ for any matrix $A$;
    \item (symmetry) $A \sim B$ if and only if $B \sim A$;
    \item (transitivity) If $A \sim B$ and $B \sim C$ then $A \sim C$.
    \end{itemize}
  \item
    \textit{Row operations preserve solutions.}  Suppose $A$ and $B$
    are the augmented matrices of two linear systems.  If $A \sim B$
    then $\vs \in \R^n$ is a solution of the linear system with
    augmented matrix $A$ if and only if it is a solution of the system
    with augmented matrix $B$.
  \end{itemize}

  \textit{Example.}
  \todo{Fill in example.}
\item
  The \textbf{leading entry} of a row in a matrix is its first nonzero
  entry.  Note that a row may not have a leading entry.

  \textit{Example.}
  \todo{Fill in example.}
\item
  A matrix is in \textbf{echelon form} if
  \begin{itemize}
  \item
    the leading entry of every row appears to the right of the
    leading entries of the rows that precede it;
  \item
    every all-zero row appears below every non-zero row.
  \end{itemize}
  This can be expressed by example pictorially as
  \begin{displaymath}
    \begin{bmatrix}
      0&\blacksquare&*&*&*&*&*&*&*&*\\
      0&0&0&\blacksquare&*&*&*&*&*&*\\
      0&0&0&0&\blacksquare&*&*&*&*&*\\
      0&0&0&0&0&\blacksquare&*&*&*&*\\
      0&0&0&0&0&0&0&0&\blacksquare&*\\
      0&0&0&0&0&0&0&0&0&0\\
    \end{bmatrix}
  \end{displaymath}
  where ``$\blacksquare$'' represents a nonzero entry and ``$*$" represents
  an entry of any value.

  \textit{Facts.}
  \begin{itemize}
  \item
    \textit{Echelon forms can be used to determine consistency.}
    Let $A$ be the augmented matrix of a linear system.  The system is
    inconsistent if and only if every echelon form of $A$ has a row of
    the form
    \begin{displaymath}
      \begin{bmatrix}
        0 \dots & 0 & b
      \end{bmatrix}
    \end{displaymath}
    where $b \neq 0$.
  \end{itemize}

  \textit{Examples.}
  \begin{itemize}
  \item
    \todo{Example}
    \todo{Example}
  \end{itemize}
\item A matrix is in \textbf{row-reduced echelon form} (RREF) if it is echelon form and
  \begin{itemize}
  \item every leading entry is $1$;
  \item its leading entries are the only nonzero entries of their columns.
  \end{itemize}
  This can be expressed by example pictorially as
  \begin{displaymath}
    \begin{bmatrix}
      0&1&*&0&0&0&*&*&0&*\\
      0&0&0&1&0&0&*&*&0&*\\
      0&0&0&0&1&0&*&*&0&*\\
      0&0&0&0&0&1&*&*&0&*\\
      0&0&0&0&0&0&0&0&1&*\\
      0&0&0&0&0&0&0&0&0&0\\
    \end{bmatrix}
  \end{displaymath}
  where ``$*$" represents an entry of any value.  We will also say
  that a matrix \textit{is} an RREF to mean that it is in row-reduced
  echelon form.

  \textit{Fact.} \textit{Every matrix is row-equivalent to exactly one RREF.}

  \textit{Examples.}
  \begin{itemize}
  \item \todo{Example}
  \item \todo{Example}
  \end{itemize}

\item A \textbf{pivot position} of a matrix is a position of a leading
  entry in its RREF.  A \textbf{pivot column} is a column containing a
  pivot position.
\item A variable $x_i$ is \textbf{basic} with respect to a linear
  system if the $i$th column of its augmented matrix is a pivot
  column.  A variables is \textbf{free} if it is not basic.
\item A \textbf{general form solution}\footnote{This is also
sometimes just called a general solution.} is a collection of statements of the form
  \begin{displaymath}
    x_i = x_i \text{ is free} \qquad \text{or} \qquad x_i = a_1x_{j_1} + \dots + a_kx_{j_k}
  \end{displaymath}
  where $a_{1}, \dots, a_{k} \in \R$ and $x_{j_1}, \dots, x_{j_k}$ are
  free variables.  In other words, a general form solution is one that
  expresses every basic variable of a linear system linearly in terms
  of its free variables.
\end{enumerate}

\subsection*{Topics Not Covered}

\begin{itemize}
\item Gaussian elimination
\item deriving general form solutions from RREFs
\end{itemize}
